% !TEX root = chapter-standalone.tex


\section{Products and sums of monotone functions}
\label{sec:monotone-products-sums}

In \cref{chap:poset-build} we constructed the product and the sum of two \SY{preorders}, and in \cref{sec:new-functions-from-old} we constructed the product and the sum of two functions.
Putting these together, we can form products and sums of \SY{monotone functions}.
Throughout this section, $\posA$, $\posB$, $\posC$, and $\posD$ are \SY{preorders}.

\subsection{Products}

\begin{ctdefinition}[Product of monotone functions]
    \label{def:monotone-product}
    Let $\mapa \colon \posA \mto \posB$ and $\mapb \colon \posC \mto \posD$ be \SY{monotone functions}.
    Their \emph{product} is the product of functions (\cref{def:product-of-functions})
    \begin{equation}
        \defmapcomma{
            \mapa \funcprod \mapb
        }{
            \posA \Ptimes \posC
        }{
            \mto
        }{
            \posB \Ptimes \posD
        }{
            \tup{\posAel, \posCel}
        }{
            \tup{\mapa(\posAel), \mapb(\posCel)}
        }
    \end{equation}
    where the source and target carry the product preorders (\cref{def:poset-product}).
\end{ctdefinition}

\begin{lemma}
    \label{lem:monotone-product}
    The product $\mapa \funcprod \mapb$ of two monotone functions is monotone.
\end{lemma}
\begin{proof}
    Suppose $\tup{\posAel, \posCel} \posleqof{\posA \Ptimes \posC} \tup{\posAel', \posCel'}$, that is, $\posAel \posAleq \posAel'$ and $\posCel \posCleq \posCel'$.
    Since $\mapa$ and $\mapb$ are monotone, $\mapa(\posAel) \posBleq \mapa(\posAel')$ and $\mapb(\posCel) \posDleq \mapb(\posCel')$.
    By the definition of the product preorder, this says $\tup{\mapa(\posAel), \mapb(\posCel)} \posleqof{\posB \Ptimes \posD} \tup{\mapa(\posAel'), \mapb(\posCel')}$.
\end{proof}

In words: the product of two monotone functions acts on each component separately, and since the order on a product is also defined component by component, monotonicity is preserved.

\begin{example}[Converting units]
    \label{exa:monotone-product-units}
    The functions converting temperatures from degrees Celsius to degrees Fahrenheit, and masses from kilograms to pounds,
    \begin{equation}
        \defmapcomma{
            \mapa
        }{
            \realswithleq
        }{
            \mto
        }{
            \realswithleq
        }{
            \ela
        }{
            \tfrac{9}{5} \ela + 32
        }
        \qquad
        \defmapcomma{
            \mapb
        }{
            \tup{\nonNegReals, \leq}
        }{
            \mto
        }{
            \tup{\nonNegReals, \leq}
        }{
            \elb
        }{
            2.2046 \, \elb
        }
    \end{equation}
    are monotone.
    Their product $\mapa \funcprod \mapb$ converts a pair $\tup{\text{temperature}, \text{mass}}$ from metric to imperial units, for example $(\mapa \funcprod \mapb)(\tup{20, 3}) = \tup{68, 6.6138}$.
    By \cref{lem:monotone-product}, if one pair is below another in the product order (warmer \emph{and} heavier), then so are the converted pairs.
\end{example}

A product of preorders comes with two \emph{projections}
\begin{equation}
    \defmapcomma{
        \proj_1
    }{
        \posA \Ptimes \posB
    }{
        \mto
    }{
        \posA
    }{
        \tup{\posAel, \posBel}
    }{
        \posAel
    }
    \qquad
    \defmapcomma{
        \proj_2
    }{
        \posA \Ptimes \posB
    }{
        \mto
    }{
        \posB
    }{
        \tup{\posAel, \posBel}
    }{
        \posBel
    }
\end{equation}
which are monotone, directly by the definition of the product preorder.
Conversely, two monotone functions with a common source can be combined into one monotone function into the product.

\begin{lemma}
    \label{lem:monotone-pairing}
    Let $\mapa \colon \posC \mto \posA$ and $\mapb \colon \posC \mto \posB$ be \SY{monotone functions}.
    Then their \emph{pairing}
    \begin{equation}
        \defmapcomma{
            \tup{\mapa, \mapb}
        }{
            \posC
        }{
            \mto
        }{
            \posA \Ptimes \posB
        }{
            \posCel
        }{
            \tup{\mapa(\posCel), \mapb(\posCel)}
        }
    \end{equation}
    is monotone, and it satisfies $\tup{\mapa, \mapb} \mthen \proj_1 = \mapa$ and $\tup{\mapa, \mapb} \mthen \proj_2 = \mapb$.
\end{lemma}
\begin{proof}
    If $\posCel \posCleq \posCel'$, then $\mapa(\posCel) \posAleq \mapa(\posCel')$ and $\mapb(\posCel) \posBleq \mapb(\posCel')$, which says $\tup{\mapa, \mapb}(\posCel) \posleqof{\posA \Ptimes \posB} \tup{\mapa, \mapb}(\posCel')$.
    The two equations hold because $\proj_1$ and $\proj_2$ pick out the first and the second component.
\end{proof}

\begin{example}[Energy and time of a trip]
    \label{exa:monotone-pairing-trip}
    An electric vehicle that drives a distance $\ela$ (in km) needs the energy $\mapa(\ela) = 0.2 \, \ela$ (in kWh) and, at an average speed of $\unit[50]{km/h}$, the time $\mapb(\ela) = \ela / 50$ (in hours).
    Both are monotone functions $\tup{\nonNegReals, \leq} \mto \tup{\nonNegReals, \leq}$, and their pairing
    \begin{equation}
        \tup{\mapa, \mapb} \colon \tup{\nonNegReals, \leq} \mto \tup{\nonNegReals, \leq} \Ptimes \tup{\nonNegReals, \leq},
        \qquad
        \ela \mapsto \tup{0.2 \, \ela, \ela / 50}
    \end{equation}
    is monotone: a longer trip needs both more energy and more time.
\end{example}

\subsection{Sums}

Recall that the elements of the sum $\posA \Pplus \posC$ (\cref{def:disjoint-union-of-posets}) are tagged elements $\tup{1, \posAel}$ and $\tup{2, \posCel}$, and that elements with different tags are never related.

\begin{ctdefinition}[Sum of monotone functions]
    \label{def:monotone-sum}
    Let $\mapa \colon \posA \mto \posB$ and $\mapb \colon \posC \mto \posD$ be \SY{monotone functions}.
    Their \emph{sum} is the sum of functions (\cref{def:sum-of-functions})
    \begin{equation}
        \begin{aligned}
            \mapa \funcsum \mapb \colon \posA \Pplus \posC & \mto \posB \Pplus \posD, \\
            \disunionA{\posAel}                             & \mapsto \disunionA{\mapa(\posAel)}, \\
            \disunionB{\posCel}                             & \mapsto \disunionB{\mapb(\posCel)},
        \end{aligned}
    \end{equation}
    where the source and target carry the sum preorders.
\end{ctdefinition}

\begin{lemma}
    \label{lem:monotone-sum}
    The sum $\mapa \funcsum \mapb$ of two monotone functions is monotone.
\end{lemma}
\begin{proof}
    Suppose that two elements of $\posA \Pplus \posC$ are related.
    By the definition of the sum preorder, they lie in the same summand.
    If they are $\disunionA{\posAel} \posleqof{\posA \Pplus \posC} \disunionA{\posAel'}$, then $\posAel \posAleq \posAel'$, so $\mapa(\posAel) \posBleq \mapa(\posAel')$, and hence $\disunionA{\mapa(\posAel)} \posleqof{\posB \Pplus \posD} \disunionA{\mapa(\posAel')}$.
    The case of two elements of the second summand is the same, with $\mapb$ in place of $\mapa$.
\end{proof}

\begin{example}
    \label{exa:monotone-sum}
    The functions $\mapa \colon \natswithleq \mto \natswithleq$, $\ela \mapsto \ela^2$, and $\mapb \colon \realswithleq \mto \realswithleq$, $\elb \mapsto \elb^3$, are monotone, and so is their sum
    \begin{equation}
        \mapa \funcsum \mapb \colon \natswithleq \Pplus \realswithleq \mto \natswithleq \Pplus \realswithleq,
        \qquad
        \disunionA{\ela} \mapsto \disunionA{\ela^2},
        \quad
        \disunionB{\elb} \mapsto \disunionB{\elb^3}.
    \end{equation}
    Note that squaring would \emph{not} be monotone on all integers, since $-2 \leq 1$ but $4 > 1$; this is why we take $\mapa$ to be defined on the natural numbers here.
\end{example}

A sum of preorders comes with two \emph{inclusions}
\begin{equation}
    \defmapcomma{
        \inj_1
    }{
        \posA
    }{
        \mto
    }{
        \posA \Pplus \posB
    }{
        \posAel
    }{
        \disunionA{\posAel}
    }
    \qquad
    \defmapcomma{
        \inj_2
    }{
        \posB
    }{
        \mto
    }{
        \posA \Pplus \posB
    }{
        \posBel
    }{
        \disunionB{\posBel}
    }
\end{equation}
which are monotone, directly by the definition of the sum preorder.
Conversely, two monotone functions with a common target can be combined into one monotone function out of the sum.

\begin{lemma}
    \label{lem:monotone-copairing}
    Let $\mapa \colon \posA \mto \posC$ and $\mapb \colon \posB \mto \posC$ be \SY{monotone functions}.
    Then their \emph{copairing}
    \begin{equation}
        \begin{aligned}
            [\mapa, \mapb] \colon \posA \Pplus \posB & \mto \posC, \\
            \disunionA{\posAel}                      & \mapsto \mapa(\posAel), \\
            \disunionB{\posBel}                      & \mapsto \mapb(\posBel),
        \end{aligned}
    \end{equation}
    is monotone, and it satisfies $\inj_1 \mthen [\mapa, \mapb] = \mapa$ and $\inj_2 \mthen [\mapa, \mapb] = \mapb$.
\end{lemma}
\begin{proof}
    Related elements of $\posA \Pplus \posB$ lie in the same summand.
    If $\disunionA{\posAel} \posleqof{\posA \Pplus \posB} \disunionA{\posAel'}$, then $\posAel \posAleq \posAel'$, so $[\mapa, \mapb](\disunionA{\posAel}) = \mapa(\posAel) \posCleq \mapa(\posAel') = [\mapa, \mapb](\disunionA{\posAel'})$; the second summand is treated in the same way, using $\mapb$.
    The two equations hold by definition.
\end{proof}

Note that nothing needs to be checked for pairs of elements from different summands, because these are never related.
This is what makes it possible to combine two completely unrelated monotone functions into one.

\begin{example}[Postage]
    \label{exa:monotone-copairing-postage}
    A post office sends parcels, whose price depends on their weight, and letters, whose price depends on their number of pages.
    Let $\mapa \colon \tup{\nonNegReals, \leq} \mto \tup{\nonNegReals, \leq}$ give the price of a parcel as a function of its weight, and $\mapb \colon \natswithleq \mto \tup{\nonNegReals, \leq}$ give the price of a letter as a function of its number of pages; both are monotone, since heavier parcels and longer letters cost more.
    The copairing
    \begin{equation}
        [\mapa, \mapb] \colon \tup{\nonNegReals, \leq} \Pplus \natswithleq \mto \tup{\nonNegReals, \leq}
    \end{equation}
    is the price list of the post office for all items, parcels and letters alike.
    It is monotone, although a parcel and a letter cannot be compared with each other: in the sum, they are simply not related.
\end{example}
